\documentclass[twoside,11pt]{article}

% Any additional packages needed should be included after jmlr2e.
% Note that jmlr2e.sty includes epsfig, amssymb, natbib and graphicx,
% and defines many common macros, such as 'proof' and 'example'.
%
% It also sets the bibliographystyle to plainnat; for more information on
% natbib citation styles, see the natbib documentation, a copy of which
% is archived at http://www.jmlr.org/format/natbib.pdf

\usepackage{jmlr2e}
\usepackage{mathtools}
\usepackage{amssymb}
\usepackage{bm}
\usepackage{preamble}

% todo
 \setlength {\marginparwidth }{2cm} % TODO: REMOVE THAT later
\usepackage{todonotes}
\definecolor{lightred}{RGB}{120,60,60}
\definecolor{darkred}{RGB}{255,0,0}
\newcommand{\proofread}[1]{{\color{lightred}#1}}
\newcommand{\rewrite}[2]{{\color{darkred}#1}}

% Define softer colors
\definecolor{softgreen}{RGB}{102, 163, 102}  % A muted green
\definecolor{softred}{RGB}{204, 102, 102}    % A muted red
\definecolor{softblue}{RGB}{102, 153, 204}   % A muted blue
\definecolor{softpurple}{RGB}{153, 102, 204} % A muted purple
\definecolor{softorange}{RGB}{224, 153, 102} % A softer orange
\definecolor{softgray}{RGB}{150, 150, 150}   % A medium gray

\definecolor{gbcolor}{RGB}{220, 160, 180}
\newcommand{\gb}[2][]{\todo[color=gbcolor, caption={GB}, #1]{#2}}
\definecolor{jbcolor}{RGB}{180, 220, 200}
\newcommand{\jb}[2][]{\todo[color=jbcolor, caption={JB}, #1]{#2}}
\definecolor{bpcolor}{RGB}{255, 215, 0}
\newcommand{\bp}[2][]{\todo[color=bpcolor, caption={BP}, #1]{#2}}

% Definitions of handy macros can go here

\newcommand{\dataset}{{\cal D}}
\newcommand{\fracpartial}[2]{\frac{\partial #1}{\partial  #2}}

\newcommand{\E}{{\mathbb{E}}}
\newcommand{\R}{{\mathbb{R}}}
\newcommand{\te}{{\theta}}


% Heading arguments are {volume}{year}{pages}{submitted}{published}{author-full-names}

\jmlrheading{1}{2000}{1-48}{4/00}{10/00}{Blanc}

% Short headings should be running  head and authors last names

\ShortHeadings{A Moment, General!}{Blanc, Bodelet \& Poilane }
\firstpageno{1}

\begin{document}

\title{A Generalized Method of Moments for Deep Latent Variable Models}

\author{\name Guillaume Blanc \email guillaume.e.blanc@gmail.com
  \AND
  \name Julien Bodelet \email julien.bodelet@gmail.com
  \AND
  \name Benjamin Poilane \email benjamin@protonmail.com 
}

\editor{Fan Lui Hui} % lol

\maketitle

\begin{abstract}%   <- trailing '%' for backward compatibility of .sty file
This paper introduces \gb[]{find good name}, a new framework, from theoretical justifications to practical implementations, for consistent estimation of latent variables models. Using the statistical theory of M-estimators and leveraging the universal approximation of neural networks, empowered by technological innovations in auto-diff software, we provide alternatives to variational inference and adversarial methods to estimate deep generative models. We demonstrate this methodology on a non-parametric latent variable model and the traditional MNIST dataset.
\end{abstract}

\begin{keywords}
Latent variable models
\end{keywords}

\newpage
\section*{Goals}

\begin{itemize}
    \item [Objectif 1]:convergence de theta hat avec p fixé et $n \to \infty$.
    Cette convergence est en sqrt(n) avec sandwitch asymptotic variance si on a :
    \begin{itemize}
        \item [1)] unicité de la solution des estimating equations. Questions:
        \begin{itemize}
             \item est-ce que unicité de la solution du modèle marginal (optimization de la marginal likelihood) implique unicité de solution des EE ? Sous quelles conditions?
             \item  peut-on se passer de l'unicité ? Et si on suppose qu'on a un moyen de converger vers le vrai $\theta_0$ ?
             \item  est-ce que si on suppose p suffisamment large (mais fixe), on peut avoir e(y, theta) suffisamment proche des vrais z, et en déduire unicité des solutions ? Y-a-til un lien à faire avec la robustesse ?
        \end{itemize}
        \item [2)] l'espérance de la dérivée du score non singulière. Questions :
        \begin{itemize}
            \item  Est-ce que la non singularité de l'espérance de la dérivée du score associé à l'extended likelihood (z connus) permet d'avoir 2) ? Explorer le score-trick et les conditions nécessaires sur la fonction e(., .) pour ça.
        \end{itemize}
    \end{itemize}

    \item [Objectif 2]: Si on autorise p à diverger en même temps que n. A-t-on convergence de theta hat gràce à la convergence de e(y, theta) vers z ? On pourrait rapprocher les EE de celles d'un modèle où les Z sont connus où on a des résultats connus (simple GLM).
\end{itemize}
\subsection*{Who does what}
\bp[inline]{Work on objective 1.2}
\jb[inline]{Work on objective 1.1}
\gb[inline]{Work on objective 1.2}


\newpage

\listoftodos

\newpage


\section{Introduction}
 Latent variable models are commonly used to capture complex dependencies in multivariate data. Generating a data point $\Y\in \mathcal Y \subseteq \mathbb R^p$ from a latent variable model  involves two steps: (1) a latent (unobserved) variable $\Z\in \mathcal Z \subseteq\mathbb R^q$ is generated from a known prior distribution $\pi(\z)$; (2) given $\Z$, the random variable $\Y|\Z$ is generated from some conditional distribution $f_\bt(\y|\z)$, known up to the parameters $\bt\in\bm\Theta\subseteq \mathbb R^r$\jb{Does it work for nonparametric?}. Thus, the latent variable model is a \textit{generative model} $F_\bt$ for the pair $(\Y, \Z)$:
%
\begin{align}
    \begin{matrix}
        \hfill\Z &\sim& \pi(\z)\hfill\\
        \hfill\Y|\Z &\sim& f_\bt(\y|\z)\hfill,
    \end{matrix}
    \label{eqn:generative-model-general}
\end{align}
%
where $\Z$ is unobserved.
% 
The marginal density $f_\bt(\y)$ of $\Y$ is then obtained by integrating out the latent variable, 
%
\begin{equation}
    f_\bt(\y) = \int_{\mathcal Z} f_\bt(\y|\z)\pi(\z)d\z.
\end{equation}
%
Latent variable models are elegant tools for multivariate data analysis, are easy to set up, and are surprisingly general. Identifiability usually requires that the dimension of the latent variable space, $q$, be smaller than that of the feature space, $p$. Far from being a drawback, this constraint enables dimensionality reduction, where a small number of latent variables capture the most salient features of the data. As such, they have been used extensively in deep learning, etc... In the social and biological sciences, they are a key tool to estimate a theorized latent construct ($\Z$) from a set of observable features ($\Y|\Z$) that it is assumed to affect. In longitudinal settings, where many observations of the same unit are gathered over time, a latent variable known as random effect can be used to model the correlation between observations on the same unit.




% we don't need to talk (much) about the marginal likelihood, since the method only relies on the extended likelihood.
% I would just talk about $f_\theta(y)$, which is intractable due to the integral.
Given $n$ independent data pairs $\{(\Y_i^0, \Z_i^0)\}_{i=1}^n$ from $F_{\bt_0}$, where $\Z_i^0$ are unobserved, an estimate of $\bt_0$ can, in principle, be obtained by maximizing the sample \textit{marginal log-likelihood}, 
%   
\begin{equation}
\sum_{i=1}^n l(\bt; \Y_{i}^0)= \sum_{i=1}^n \log\left(\int_{\mathcal Z}f_\bt(\Y_{i}^0|\z)\pi(\z)d\z\right),
    \label{eqn:marginal-likelihood}
\end{equation}
%
where $l(\bt; \y)$ is the marginal log-likelihood function, which we define for a single observation.
A chief challenge in estimating latent variable models is that the multiple integrals prevent the direct maximization of \eqref{eqn:marginal-likelihood}; indeed, apart in simple cases which admit analytical expressions in its derivation, even the Expectation-Maximization algorithm, despite its good properties, is computationally too intensive to solve. Alternatives include quadrature of the integral, using lower bounds on the marginal likelihood (VAE...), but these approximations are often not adapted for larger dimensions, and may include severe bias in the estimates.



\section{The Framework}


\subsection{Notation}
We write random variables in upper-case letters and function arguments in low-case letters. Bold characters reflect to vector or matrix quantities. For instance the density of a vector random variable $\Y$ may be written $f_\bt(\y)$ where $\bt$ is  a vector of parameters. 


\subsection{The Estimator}
The estimator is the root of  

$$
T(y) E_{Z\sim q(z|y;\theta)}\left[\frac{\partial \eta(z;\theta)}{\partial \theta}\right] - E_{Y\sim f(y;\theta)}\left[T(Y) E_{Z\sim q(z|Y;\theta)}\left[\frac{\partial \eta(z;\theta)}{\partial \theta}\right]\right] = 0
$$

It is:
\begin{itemize}
    \item Fisher consistent by construction
    \item robust to misspecification of $q$ (it can be a dirac centered on the map, it can be an approximation
\end{itemize}


\subsection{The EM Algorithm}

The Expectation-Maximization (EM) algorithm is an iterative method for estimating parameters in latent variable models, especially when direct maximization of the marginal log-likelihood \eqref{eqn:marginal-likelihood} is challenging due to unobserved variables.

The EM algorithm alternates between the following two steps:

\begin{itemize}
    \item \textbf{E-step:} Given the current parameter estimate \(\bt^{(t)}\), compute the conditional expectation of the complete-data log-likelihood:
    \[
    Q(\bt|\bt^{(t)}) = \mathbb{E}_{\Z|\Y, \bt^{(t)}}\left[ \log f_\bt(\Y, \Z) \right].
    \]

    \item \textbf{M-step:} Maximize \(Q(\bt|\bt^{(t)})\) with respect to \(\bt\):
    \[
    \bt^{(t+1)} = \arg\max_\bt Q(\bt|\bt^{(t)}).
    \]
\end{itemize}

At convergence, the solution satisfies the stationary condition that the derivative of \(Q(\bt|\hat{\bt})\) with respect to \(\bt\), evaluated at \(\bt = \hat{\bt}\), is zero:
\[
\left. \nabla_\bt Q(\bt|\hat{\bt}) \right|_{\bt = \hat{\bt}} = 0.
\]
%
Under suitable conditions, \(\hat{\bt}\) is a local maximizer of the marginal likelihood \eqref{eqn:marginal-likelihood}.


Computationally, a key challenge is that the EM algorithm does not directly apply due to the intractability of the complete-data log-likelihood. This arises because the conditional expectation in the E-step requires complex integrals over the latent variable \(\Z\), which lack closed-form solutions. Thus, alternative approaches, such as variational approximations or stochastic methods, are often needed to approximate the E-step, or the marginal likelihood itself.


In this work, we propose a method to


\gb[inline, caption= improve the introduction with other methods]{This approach is used in GLMM, in GLLVM, VAE, GANs, as well as for missing data, etc... Write more.}


\subsection{Deep Generalized Latent Variable Models}

A Deep Generalized  Latent Variable Model (DGLVM)\gb[]{This doesn't exist, this is the model we want ( or not? ) to introduce. Perhaps we want to consider this model as a special case of our estimation strategy. We'll see that later.} further specifies the conditional density $f_\bt(\y|\z)$, in two ways: 
%
(i) given the latent variable $\Z$, the elements $Y_{1}, \dots, Y_{p}$ of $\Y$ are independent, that is $Y_{j} \!\! \perp \!\!  Y_{j'}\vert \Z$, $\forall j\neq j'$, which allows to re-write $f_\bt(\y| \z) = \prod_{j=1}^p f_{j, \bt_j}(y_j| \z)$, with $\bt_j\subset \bt$ the parameters specific to $f_j$, the conditional density or probability mass function for response $j$.
%
(ii) For all \(j\), the individual \(Y_j\) are distributed according to a member of the exponential family of distributions, hence $f_{j, \bt_j}(y_j| \z)$ admits the usual form
%
$$
f_{j, \bt_j}(y_j | \z) = h(y_j) \exp\left( \eta(\theta_j)^\top T(y_j) - A(\theta_j) \right),
$$
%
for known functions $h$, $\eta$, $T$ and $A$.

\subsection{A Fisher Consistent Estimator}

We define  our estimator as the root of the following estimating equations %
%
\[
T(y) \, \mathbb{E}_{Z \sim q(z \mid Y = y; \theta)} \left[ \frac{\partial \eta(Z, \theta)}{\partial \theta} \right] - \mathbb{E}_{Y \sim f(y; \theta)} \left[ T(Y) \cdot \mathbb{E}_{Z \sim q(z \mid Y; \theta)} \left[ \frac{\partial \eta(Z, \theta)}{\partial \theta} \Big| Y \right] \right] = 0
\]
%
Some properties:

\begin{itemize}
    \item Fisher consistent by definition
    \item Robust to misspecifications of q
    \item Under true q, asymptotically equivalent to MLE (?)
\end{itemize}

We optimize it using Robbins-monro stochastic approximation root-finding algorithm. The quantities can be obtained effortlessly using autodiff (as in pytorch), as we will see.

Convergence of algorithm??!! can we choose $q()$ to help? is there 

\subsubsection{equivalence mle?}
The MLE is the root of
%
\[
T(y) \, \mathbb{E}_{Z \sim q(z \mid Y = y; \theta)} \left[ \frac{\partial \eta(Z, \theta)}{\partial \theta} \right] - \mathbb{E}_{Z \sim q(z \mid y; \theta)} \left[ E[Y|\eta(Z, \theta)] \frac{\partial \eta(Z, \theta)}{\partial \theta} \right] = 0
\]
%



\subsection{Older version: extended likelihood approach}
In order to avoid the computation of the integral, an alternative way is to consider performing simultaneous inference on $\{\z_{i}\}_{i=1}^n$ and $\bt$ treating the latent variables as parameters, by maximizing the so-called sample \textit{extended log-likelihood}
%
\begin{equation}
    \sum_{i=1}^n l(\bt, \z_{i}; \Y_{i}^0) = \sum_{i=1}^n\log(f_\bt(\Y_{i}^0|\z_{i})\pi(\z_{i})),
    \label{eqn:extended-likelihood}
\end{equation}
%
where $l(\bt, \z; \y)$ is the extended log-likelihood function for a single observation. Compared with $l(\bt; \y)$, $l(\bt, \z; \y)$ does away with the integrals altogether, which makes it tractable. 
The (marginal) scores are defined by
$$
\bm\psi_n(\bt, \z) = \frac{1}{n}\sum_{i=1}^n
\bm\psi(\bt, \z_i, \y_i)
= \frac{1}{n}\sum_{i=1}^n \frac{\partial}{\partial\theta} l(\bt, \z_{i}; \Y_{i}). 
$$
In the case that  $\z$ is known, $\psi_n$ corresponds to the score of a GLM, and is Fisher consistent conditionally on $\z$:
$$
\E[\bm\psi_n(\bt, \z, \y)| \z] = 0\,.
$$
The profile likelihood treats $\z$ as a nuisance parameter. First, one obtains an estimator of $z$ given the parameter 
$\bt$ by solving the score equation:
$$
z_\bt: \bm\phi_n(\bt, \z) = 0.
$$ \bp{a hat on $z$ ?}
Then one solves the score by plugin the estimator of the nuisance parameter:
$$
\hat \bt_{\text{PL}}: \bm\psi_n(\bt, \z_\bt, \y) = 0.
$$
When it exists and $\bm\phi_n(\bt, \z) = \frac{1}{n}\sum_{i=1}^n \frac{\partial}{\partial\z} l(\bt, \z_{i}; \Y_{i})$,
the profile likelihood estimator is equivalent to the Maximum likelihood estimator.
However, while commonly used for mixed effect models, 
the profile likelihood fails for most latent variable models, since it is not Fisher consistent, i.e. $\E \bm\psi_n(\theta, \z_\theta, \y)\neq 0$.
Note that if $\z_i$ were observed, maximizing \eqref{eqn:extended-likelihood} with respect to $\bt$ would yield the maximum likelihood estimator of $\bt_0$. What's more, the extended likelihood \eqref{eqn:extended-likelihood} is often used to estimate the unknown $\Z_i^0$, obtained as the maximizer of \eqref{eqn:extended-likelihood} with respect to $\z_i$, holding $\bt$ at some estimate of $\bt_0$. 
However, unless stringent criteria are met\gb{regarding the independence of the estimates of $\Z$ and $\bt$}, simultaneously maximizing \eqref{eqn:extended-likelihood} with respect to both $\bt$ and $\{\z_{i}\}_{i=1}^n$ is not viable. 
\jb{Talk about random effect model (doable but biased).}
In simple models, a so-called \textit{canonical} scale for the latent variables may exist, aiming to make the estimates of $\bt_0$ and $\Z_i^0$ independent and enabling their simultaneous estimation using \eqref{eqn:extended-likelihood}, but this scale may be hard to define and may not even exist in more complex models \gb{cite the extended likelihood approach that defines which criteria are considered}. Other methods may include regularization of \eqref{eqn:extended-likelihood},  but the regularization parameter introduce bias in the estimates and need to be tuned, complicating the analysis \gb{cite Hastie}.

%In the following, we solve these challenges faced by $l(\bt, \z;\y)$ using the approximation-imputation (AI) trick, which is an alternative to simultaneous maximization of $\bt$ and $\z_i$ that keeps the good properties and ease of computation of the extended likelihoof function, while providing guarantees of consistent estimation.


\subsection{Partial profiling}

We propose a novel method which is based on partial profiling.

First, we assume that the conditional densities $f_{\bt, j}(\y|\z)$ belong to the exponential family of distributions, that is

\begin{equation}
   f_{\bt, j}(y|\z) = \exp(\{y\gamma - b_j(\gamma))\}/a_j(\phi_j) + c_j(y, \gamma)),
   \label{eqn:conditional-distribution}
\end{equation}
for some known functions $\{a_j, b_j, c_j\}_{j=1}^p$, arbitrary scale parameters $\{\phi_j\}_j^p$ and where $\gamma$ is called the cannonical parameter, which in turn entirely depends on $\bt$ and $\z$. 


\gb[inline]{Write more about this framework, how general it remains, etc etc. and how widely it is, in fact, used. Then introduce the maximum likelihood principle}

%\gb[inline]{Should we write $(\Y^{(i)}, \Z^{(i)})$ instead of $(\Y_i, \Z_i)$, when referring to the observations from $F_{\bt_0}$, to avoid the confusion with the $j$th element of the random variable $\Y$, i.e. $Y_j$? The only difference is the context in which it is used, and the bold letter for the observations from $F_{\bt_0}$. It can be useful to help the reader here. Alternatively, we could write $(\Y_i^0, \Z_i^0)$ to remind the reader it is from $F_{\bt_0}$.}


%We are first interested in the estimating function for $\bt$ when $\z$ is known. 
The estimating function for $\bt$ is obtained by taking the partial derivative of the extended likelihood function $l(\bt, \z; \y)$ with respect to $\bt$. This (partial) derivative takes the form of the sum of $p$ familiar GLM scores (assuming $\z$ to be known):
%
\begin{equation}
\bm\Psi_n(\bt, \z) = \frac{\partial l(\bt, \z;\y)}{\partial \bt}=\sum_{j=1}^p \bm D_j(\bt, \z)\frac{(y_{j} - \mu_j(\bt, \z))}{V_j(\bt, \z)\phi_j},
    \label{eqn:generalized-linear-models}
\end{equation}
%
%\gb[inline]{Since $\z$ is unknown, our strategy is to impute them by $\e(\y, \bt)$. However, doing so means the estimator is not fisher consistent -- that is, the estimating equations are not 0 in expectation. That's bad. Very bad. Our solution is to brute-force this by substracting their expectation. Yay, clever girl! Is that good? No, not good. Not always identifiable. So, still bad. Sad, very sad. But in fact, we can just realize after three years of trying out stuff that the mu is itself an expectation, so we can separate the estimating equation between a quantity and its expectation under the model, impute it, and VOILA! IT WORKS! NICE! REKABOOMBOOM}
%
where, for all $j$, $\mu_j(\bt, \z) = \E[Y_{j}|\z] = g_j^{-1}(\eta_j(\bt, \z))$, $\bm D_j(\bt, \z) = \partial \mu_j(\bt, \z)/\partial \bt$, $V_j(\bt, \z)\phi_j = \text{Var}(Y_{j}|\z)$, $g_j$ is a link function, and $\phi_j$ is an additional variance parameter, and where the conditional expectation and variance are taken with respet to the conditional density $f_{\bt,j}(y|\z)$.



Consider the score equations for the exponential framework:
\begin{equation}
\bm\Psi_n(\bt, \z) = \sum_{j=1}^p\frac{\bm D_j(\z, \bt) (y_{ij}- \mu_j(\bt, \z_i))}{V_j(\z_i, \bt)\phi_j}.
\label{eqn:Psi}
\end{equation}
where, for all $j$, $\mu_j(\bt, \z) = \E[Y_{j}|\z] = g_j^{-1}(\eta_j(\bt, \z))$, $\bm D_j(\bt, \z) = \partial \mu_j(\bt, \z)/\partial \bt$, $V_j(\bt, \z)\phi_j = \text{Var}(Y_{j}|\z)$, $g_j$ is a link function, and $\phi_j$ is an additional variance\bp{overdispersion parameter ? GB: pas nécessairement: il est tel que $V_j(\bt,\z)\phi_j = Var(Y_j|\z)$: il doit être estimé aussi... mais avec d'autres estimating equations. faut traiter ca.} parameter, and where the conditional expectation and variance are taken with respet to the conditional density $f_{\bt,j}(y|\z)$.
For brevity we write
\begin{equation}
\bm\Psi_n(\bt, \z) = \frac{1}{n} \sum_{i=1}^n\sum_{j=1}^p \bm W_j(\z_i, \bt) (y_{ij}- \E_\bt(y_{ij}|z_i)).
\label{eqn:score-known-z-simple}
\end{equation}
where we replaced $\mu_j(\bt, \z)$ by $\E_\bt(y|z)$.
We then note that:
\begin{equation}
\bm\Psi_n(\bt, \z) = \frac{1}{n} \sum_{i=1}^n\sum_{j=1}^p \bm W_j(\z_i, \bt) y_{ij}- \E_\bt( \bm W_j(\z_i, \bt) y_{ij}|z_i ).
\label{eqn:score-known-z-simple_2}
\end{equation}
We then compute the partial profiling, i.e. we replace the $\z_\bt$ only in the weights in the score:
\begin{equation}
\tilde{\bm\Psi}_n(\bt, \z) = \frac{1}{n} \sum_{i=1}^n\sum_{j=1}^p \bm W_j(\z_{\bt,i}, \bt) y_{ij}- \E_\bt(\bm W_j(\z_{\bt,i}, \bt) y_{ij}|\z_i) ).
\end{equation}
In contrast to full profiling, the score function
$\tilde{\bm\Psi}_n(\bt, \z)$ still depend on (the true) $\z$ through the conditional expectation.
However we argue that the dependence vanishes as $n$ diverge.
Indeed, we have 
\begin{align}
\tilde{\bm\Psi}_n(\bt, \z) 
&=
\frac{1}{n} \sum_{i=1}^n\sum_{j=1}^p \bm W_j(\z_{\bt,i}, \bt) y_{ij} - 
\frac{1}{n} \sum_{i=1}^n\sum_{j=1}^p 
\E_\bt(\bm W_j(\z_{\bt,i}, \bt) y_{ij}|\z_i),\nonumber\\    
&=
\frac{1}{n} \sum_{i=1}^n\sum_{j=1}^p \bm W_j(\z_{\bt,i}, \bt) y_{ij} - 
\E_\bt \sum_{j=1}^p \bm W_j(\z_{\bt,i}, \bt) y_{ij} + \bigO_p\left(\frac{1}{\sqrt{n}}\right)\\
&= \tilde{\bm\Psi}_n(\bt) +\bigO_p\left(\frac{1}{\sqrt{n}}\right).\nonumber
\end{align}
Basically, optimizing $\tilde{\bm\Psi}_n(\bt) $ is asymptotically equivalent to optimize $\tilde{\bm\Psi}_n(\bt, \z) $.
We thus use $\tilde{\bm\Psi}_n(\bt)$ as estimating equation for $\bt$.
We consider below the example of factor analysis which leads to explicite equations and a very fast algorithm.
Note that in general $\E_\bt$ cannot be computed analytically.
We consider later the general scenario.



\section{Generalization and connections}

\subsection{Deriving the Estimating Function}

 Writing

\begin{equation}
    \bm U(\bt, \z;\y) = \sum_{j=1}^p\frac{\bm D_j(\bt, \z) y_{j} }{V_j(\bt, \z)\phi_j},
    \label{eqn:Psi}
\end{equation}
%
the estimating function \eqref{eqn:generalized-linear-models} can be re-written WLOG as:

%
\begin{equation}
\bm s(\bt, \z; \y) = \bm U(\bt, \z; \y) - \E\left[\bm U(\bt, \z; \Y)\Large\vert\z \right].
\label{eqn:estimating-functions-known-z}
\end{equation}
%
This representation highlights a fundamental property of this estimating function, when applied to the data $(\Y, \Z)$ sampled from $F_\bt$: it is $\bm 0$ in expectation, i.e. it satisfies
%
$$
\E[\bm s(\bt, \Z; \Y)|\Z] = \bm 0,
$$
%
where the expectation is taken with respect to $f_\bt(\y|\z)$; that is, the estimating function \eqref{eqn:estimating-functions-known-z} satisfies the first Bartlett identity, one of the fundamental properties for consistent estimation. If the pairs $(\Y_i^0, \Z_i^0)$ were fully observed for all $i$, the maximum likelihood estimator for $\bt_0$, assuming it exists, would be obtained as the root in $\bt$ of
%
$\sum_{i=1}^n \bm s(\bt, \Z_i^0; \Y_i^0)$.
%

 We want to generalize the estimating function \eqref{eqn:estimating-functions-known-z} to the case where the random variable $\Z$ is not observed: it then needs to be marginalized out using its conditional distribution $f_\bt(\z|\y)$. The estimating function then becomes
%
\begin{equation}
\bm s(\bt; \y) = \int_{\mathcal Z}\bm s(\bt, \z; \y)f_\bt(\z|\y)d\z.
\label{eqn:score_function_marginal}
\end{equation}
%
It is easy to see that when applied to the data $\Y \sim f_\bt(\y)$, this marginal estimating function is still $\bm 0$ in expectation (now on the marginal model):
%
\begin{multline}
\E[\bm s(\bt; \Y)] = \int_{\mathcal Y}\int_{\mathcal Z}\bm s(\bt, \z; \y)f_\bt(\z|\y)d\z f_\bt(\y)d\y 
\\= \int_{\mathcal Z}\int_{\mathcal Y}\bm s(\bt, \z; \y)f_\bt(\y|\z)d\y \pi(\z) d\z = 0,
\end{multline}
%
where the exchange of the order of integration is justified by Fubini's theorem\gb[]{Make sure it's ok to do so}. Using \eqref{eqn:score_function_marginal} present tremendous difficulties due to the conditional expectation which is almost never tractable. A way forward would be to use variational inference, or other approximations of $\bm s(\bt, \z; \y)$. Here, we propose the following to approximate \eqref{eqn:score_function_marginal}:
%
\begin{equation}
    \bm\Psi(\bt;\y) = \bm U(\bt, \e(\bt, \y); \y) - \E[\bm U(\bt, \e(\bt, \Y); \Y)],
    \label{eqn:estimating-equations-AI}
\end{equation}
%
where $\e(\bt, \Y)$ is an estimator of $\Z$. Unlike approximations using variational inference or gaussian quadrature, the (marginal) expectation of \eqref{eqn:estimating-equations-AI} remains $\bm 0$ when applied to $\Y \sim f_\bt(\y)$:
%
\begin{equation}
    \E[\bm\Psi(\bt;\Y)] = \bm 0.
\end{equation}
%
Our strategy thus involves substituting in $\bm U$ the unknown $\Z$ by $\e(\bt, \Y)$. Finally, we define our estimator of $\bt_0$ to be the root in $\bt$ of $\sum_{i=0}^n \bm\Psi(\bt;\Y_i^0)$, that is our estimator satisfies
%
\begin{equation}
    \hat\bt = \underset{\bt\in\bm\Theta}{\text{argzero}} \sum_{i=1}^n \bm \Psi(\bt; \Y_i^0).
    \label{eqn:estimator}
\end{equation}
%

\gb[inline]{Is it possible to quantify the approximation error somehow?}


\subsection{Derivative of Estimating Equations}
We have that

\begin{align}
    \frac{\partial\bm\Psi(\bt;\y)}{\partial\bt^\top} &= \frac{\partial \bm U(\bt, \e(\bt, \y), \y)}{\partial \bt^\top} - \frac{\partial \E_\bt\left[\bm U(\bt, \e(\bt, \Y), \Y)\right]}{\partial \bt^\top}\\
    &= \frac{\partial \bm U(\bt, \e(\bt, \y), \y)}{\partial \bt^\top} - \E_\bt\left[\frac{\partial \bm U(\bt, \e(\bt, \y), \y)}{\partial \bt^\top}\right] - \E_\bt\left[\bm U(\bt, \e(\bt, \Y), \Y) \bm \psi(\bt, \Y)^\top\right]
    \label{eqn:dpsi}
\end{align}

For asymptotic normality, we need that 
\begin{align}
    -\E_{\bt}\left[\frac{\partial\bm\Psi(\bt;\Y)}{\partial\bt^\top}\right] = \E_\bt\left[\bm U(\bt, \e(\bt, \Y), \Y) \bm \psi(\bt, \Y)^\top\right]
    \label{eqn:dpsi}
\end{align}

exist and be nonsingular. The problem is that this freaking score $\bm\psi$ is an absolute pain in the ass to deal with. If only we could use the conditional scores instead.

\gb[inline]{\textbf{An argument}

Define 
\begin{equation}
    \bm A_n(\bt_0, \{\Y_i^0\}_{i=1}^n) = \frac{1}{n}\sum_{i=1}^n \E\left[-\frac{\partial}{\partial\bt^\top}\bm\Psi(\bt_0;\Y_i^0)\right]
    \label{eqn:An}
\end{equation}

\begin{equation}
    \label{eqn:An_Z}
\end{equation}



If we can have conditions that imply

\begin{equation}
    \frac{1}{n}\sum_{i=1}^n\E_{\Y\sim f_\bt(\y|\Z)}\left[\bm U(\bt, \e(\bt, \Y), \Y)|\Z\right] \to \E_{\Y\sim f_\bt(\y)}\left[\bm U(\bt, \e(\bt, \Y), \Y)\right] =\vcentcolon \bm U(\bt)
\end{equation}

and 

\begin{equation}
    \frac{1}{n}\sum_{i=1}^n\frac{\partial}{\partial \bt^\top}\E_{\Y\sim f_\bt(\y|\Z)}\left[\bm U(\bt, \e(\bt, \Y), \Y)|\Z\right] \to \frac{\partial \bm U(\bt)}{\bt^\top}
\end{equation}

then this would help, because

\begin{multline}
    \frac{\partial}{\partial \bt^\top}\E_{\Y\sim f_\bt(\y|\Z)}[\bm U(\bt, \e(\bt, \Y), \Y)|\Z] =\\ \E_{\Y\sim f_\bt(\y|\Z)}[\frac{\partial}{\partial \bt^\top}\bm U(\bt, \e(\bt, \Y), \Y)|\Z] \\ + \E_{\Y\sim f_\bt(\y|\Z)}[\bm U(\bt, \e(\bt, \Y), \Y) s(\bt, \Z;\Y)^\top|\Z] 
\end{multline}
where the score function $s(\bt, \z; \y)$ is that of the conditional model (the derivative of the log of the extended likelihood), which is easy to deal with. Maybe this helps in finding conditions on $\e$ that makes it nonsingular. 

We're interested in the sum
\begin{align}
    \frac{1}{n}\sum_{i=1}^n \bm\Psi(\bt;\Y_i) = \frac{1}{n}\sum_{i=1}^n \bm U(\bt, \e(\bt, \Y_i); \Y_i) - \frac{1}{n}\sum_{i=1}^n \E[\bm U(\bt, \e(\bt, \Y); \Y)]\\
    =\frac{1}{n}\sum_{i=1}^n \bm U(\bt, \e(\bt, \Y_i); \Y_i) - \frac{1}{n}\sum_{i=1}^n \E[\E[\bm U(\bt, \e(\bt, \Y); \Y)|\Z]]
    \label{eqn:estimating-equations-AI}
\end{align}

}

$$ \E[\frac{\partial}{\partial \theta}\bm \Psi(\theta; Y) ] = \E[\E[\frac{\partial}{\partial \theta}\bm \Psi(\theta; Y)\mid Z ]]$$



So : 

If $$ \E[\frac{\partial}{\partial \theta}\bm U(\theta, Z; Y) s^T(\theta, Z; Y)]$$ is non-singular, then $$ \E[\frac{\partial}{\partial \theta}\bm U(\theta; Y) ]$$ is non singular, because the 2 quantities are equal.


\subsection{Presentation with Marginalization}
Since $\z$ are unknown, we consider to approximate the scores $\bm s$ by
%
\begin{equation}
 \bm s(\bt; \y) = \E_{\Z\sim q_\bt(\z|\y)}\left[\bm s(\bt, \Z; \y)\right],
\end{equation}
%
where $q_\bt(\z|\y)$ defines a propositional distribution that approximates the true posterior distribution $f_\bt(\z|\y)$. A fundamental problem of $\bm s(\bt; \y)$ is that unless $q_\bt(\z|\y) = f_\bt(\z|\y)$, where the latter is the true posterior, it is not $\bm 0$ in expectation, and therefore does not define a Fisher-consistent estimator. Let's correct this!

The above can be re-written:
\begin{equation}
 \bm s(\bt; \y) = \E_{\Z\sim q_\bt(\z|\y)}\left[\bm U(\bt, \Z;\y)\right] - \E_{\Z\sim q_\bt(\z|\Y)}\left[\E_{\Y\sim f_\bt(\y|\Z)}\left[U(\bt, \Z;\Y)\right]\right]
\end{equation}
%
which can be approximated by
%
\begin{equation}
 \bm \Psi(\bt; \y) = \E_{\Z\sim q_\bt(\z|\y)}\left[\bm U(\bt, \Z;\y)\right] - \E_{\Y\sim f_\bt(\y)}\left[\E_{\Z\sim q_\bt(\z|\Y)}\left[U(\bt, \Z;\Y)\right]\right],
\end{equation}
%
which is $\bm 0 $ by definition.


\subsection{The EM Algorithm}
The EM algorithm proceeds in two steps: for a given $\bt^t$, the next iterate $\bt^{t+1}$ is obtained as follows:

\begin{itemize}
    \item [E-step]: define
    \begin{equation}
        Q(\bt | \bt^t) = \sum_{i=1}^n E_{\Z\sim f_{\bt^{t}}(\z|\Y_i^0)}[l(\bt, \Z_i;\Y_i^0)],
        \label{eqn:E-step}
    \end{equation}
    \item [M-step]: maximize:
    \begin{equation}
        \bt^{t+1} = \underset{\bt}{\text{ argmax }} Q(\bt|\bt^t)
        \label{eqn:M-step}
    \end{equation}
\end{itemize}

\subsection{Presentation As an Algorithm}


\begin{itemize}
    \item [A-step]: Approximate the gradient:
        $$
        \bm \Psi(\bt, \z; \y) = \bm U(\bt, \z; \y) - \E_{\Z\sim \pi(\z)}\left[\E_{\Y\sim f_\bt(\y|\Z)}\left[\bm U(\bt, \Z;\Y)|\Z\right]\right]
        $$
    \item [I-step]: Given current estimate $\bt^t$, impute the gradient:
        $$
        \bm\Psi(\bt, \bm \Psi(\bt, \z; \y) = \bm U(\bt, \z; \y) - \E_{\Z\sim \pi(\z)}\left[\E_{\Y\sim f_\bt(\y|\Z)}\left[\bm U(\bt, \Z;\Y)|\Z\right]\right]
        $$
    
\end{itemize}




\subsection{The Approximation-Imputation Trick}
\gb[inline]{
WARNING: The following presentation of the A-I trick may be a bit confusing. We keep it as reference, but you may not want to read it to maintain sanity.
}
In practice, the latent variables $\Z_i^0$ are not known, and $\bm s(\bt, \Z_i^0, \Y_i^0)$ cannot be computed, for two reasons: \textit{i)} the expectation is taken conditionally on the unknown $\Z_i^0$,  and \textit{ii)} the $\bm U$ function depends itself on $\Z_i^0$. We solve these problems using the Approximation-Imputation (AI) trick:




\begin{itemize}
    \item[A-step]: In the A-step (approximation), we leverage the knowledge of the (marginal) distribution of $\Z$ to approximate \eqref{eqn:estimating-functions-known-z} by 
    %
    \begin{equation}
    \bm U(\y, \z, \bt) -\E\left[\E\left[\bm U(\Y, \Z, \bt)\big|\Z \right]\right],
    \label{eqn:estimating-functions-known-z-approximated}
    \end{equation}
    %
    removing the dependence of the expectation on $\z$.
    \jb[inline]{
    This is justified by the fact that 
    \begin{equation}
        \frac{1}{n}\sum_{i=1}^n \left(\E\left[\bm U(\Y, \Z_i, \bt)\big|\Z_i \right] - \E\left[\E\left[\bm U(\Y, \Z, \bt)\big|\Z \right]\right]\right) \to \bm 0.
    \end{equation}
    (e.g. we can use the central limit theorem). In fact, using the CLT, we can actually quantify the error of this approximation, maybe that can be useful.
    }
    \item[I-step]:
    In the I-step (imputation), we impute the quantity $\z$ in $\bm U$ by an estimate, written as $\z^\ast = \bm e(\y, \bt)$, where $\e: \mathcal Y \times \mathbb R^r \to \mathbb R^q$ is a function of the data and the parameter which we call the \textit{encoder}.
    \jb[inline]{Here, it would be useful to consider typical convergence rates of $e(\Y, \bt)$ as an estimator of $\Z$. For linear models (GLLVMs) that's easy, but for more complex models, we should check. We could have an argument of the like "convergence as if we knew the true value of $\Z$", similarly to our paper. }
\end{itemize}

\noindent Using the AI trick on \eqref{eqn:estimating-functions-known-z} yields the following estimating function for the parameter $\bt$:
%
\begin{equation}
\bm \Psi(\y, \bt) = \bm U(\y, \bm e(\y, \bt), \bt) - \E\left[\bm U(\Y, \bm e(\Y, \bt), \bt)\right],
\label{eqn:estimating-function-AI}
\end{equation}
%
where the expectation is taken with respect to the marginal density $f_\bt(\y)$, making it trivial to estimate using Monte Carlo simulations due to the ease of generating a data-point from $F_\bt$. Notice that the AI-trick effectively removes the dependence on the unknown 
$\z$. Finally, we define our estimator of $\bt_0$ to be the root of the empirical mean of this function of the data, that is,
%
\begin{equation}
    \hat\bt = \underset{\bt\in\bm\Theta}{\text{argzero}} \sum_{i=1}^n \bm \Psi(\Y_i^0, \bt).
    \label{eqn:estimator}
\end{equation}
%

Note that by construction, $\E\left[ \bm \Psi( \Y, \bt)\right] = 0$, which is the key property to insure (Fisher) consistency of $\hat\bt$. What's more, the estimator \eqref{eqn:estimator} is a  This was indeed the key idea in my thesis \gb{cite thesis}.

\jb[inline]{Here we can have the asymptotic distribution using the sandwich formula for small $p$. For large $p$, using the AI-trick, we could derive other convergence properties. Whatcha think?}

\gb[inline, caption=Convergence of Approximation]{It may appear surprising that the expectation in  \eqref{eqn:estimating-functions-known-z-approximated}, which acts as a centering quantity, does not depend on $i$ -- and thus on the observation. In Theorem 1, we show that the approximation is in fact tight as $n\to\infty$.}


\subsection{Cannonical Link Functions}
In the case of a canonical link function and if $V_j$ is the standard variance function of the assumed conditional distribution for feature $j$, the following equality holds
%
$$
\bm D_j(\z, \bt) = V_j(\z, \bt)\frac{\partial \eta_j(\z, \bt)}{\partial \bt},
$$
%
and \eqref{eqn:Psi} simplifies into
%
\begin{equation}
     \sum_{j=1}^p\frac{\partial \eta_j(\z, \bt)}{\partial \bt}\frac{y_{j} }{\phi_j},
    \label{eqn:Psi-cannonical}
\end{equation}
%
where, again, $\eta_j(\z, \bt) = g_j(\mu_j(\z, \bt))$. This simplifies the computations, as we do not need to evaluate the quantities $V_j(\z, \bt)$ and $\partial \mu_j(\z, \bt)/\partial \eta_j$.

\section{Examples}
\subsection{Factor Analysis}
\gb[inline]{For now, this section is work in progress -- please improve and comment.}

Factor analysis (FA) defines the following generative model for the pair $(\Y, \Z)$:

\begin{align}
    \begin{matrix}
        \hfill\Z &\sim& \bm N(\bm 0, \bm I_q)\hfill\\
        \hfill\Y|\Z &\sim& \bm N(\bm\Lambda\Z, \bm\phi)\hfill,
    \end{matrix}
    \label{eqn:generative-model-general}
\end{align}



In factor analysis (FA) \gb{Write the generative model} the estimating equations for FA for $\bt$ can be written as
%
$$
\bm s(\bt, \z; \y) = \bm\phi^{-1}(\y - \bm\Lambda\z)\z^\top
$$
%
where $\bm\Lambda$ are the loadings and $\bm\phi$ is the diagonal matrix with diagonal elements $(\phi_1, \dots, \phi_p)$. We recognize that 
%
$$
\bm U(\y, \z, \bt) = \bm\phi^{-1}\y\z^\top
$$
%
We can consider imputing the latent variables by their maximum a posteriori (MAP), given by: 
$$
\e(\y, \bt) = (\bm\Lambda \bm\phi^{-1}\bm\Lambda^\top + \bm I_q)^{-1}\bm\phi^{-1}\bm\Lambda^\top\y =\vcentcolon  \bm K \y
$$
And thus, the estimating function is

\begin{multline}
\bm \Psi(\y, \bt) = \bm U(\y, \e(\y, \bt), \bt) - \E[\bm U(\Y, \e(\Y, \bt), \bt)] 
\\= \bm\phi^{(-1)}\y\y^\top \bm K^\top - \E\left[\bm\phi^{(-1)}\y\y^\top \bm K^\top\right]
\\= \bm\phi^{(-1)}\y\y^\top \bm K^\top - \bm\phi^{(-1)}\bm\Sigma \bm K^\top,
\end{multline}
where $\E[\Y\Y^\top] = \bm\Lambda \bm\Lambda^\top + \bm\phi = \bm \Sigma$. This estimating function is the same as the score function for the MLE of the marginal model.

Furthermore, we can now show that imputing the conditional scores $\bm s$ directly does not work:
$$
\bm s(\bt, \e(\y, \bt); \y) = (\y - \bm\Lambda \bm K \y)\y^\top \bm K^\top = \y\y^\top \bm K^\top  - \bm\Lambda \bm K \y \y^\top \bm K^\top 
$$

We now show that the imputed conditional scores are not 0 in expectation. We can use the following result easily shown using Woodburry's matrix identities:

$$
\bm \Sigma^{-1} \bm\Lambda = \bm K^\top \Leftrightarrow \bm\Lambda = \bm \Sigma \bm K^\top
$$
which is justified since $\bm\Sigma$ is invertible. Replacing $\bm\Lambda$ with this quantity, the estimating functions become

$$
 \y\y^\top \bm K^\top  - \bm \Sigma \bm K^\top \bm K \y \y^\top \bm K^\top 
 $$

 with expectation
$$
 \bm\Sigma \bm K^\top  - \bm \Sigma \bm K^\top \bm K \bm\Sigma \bm K^\top 
 $$

Using the Woodburry equality again: $\bm K^\top \bm K = \bm\Sigma^{-1} \bm\Lambda\bm\Lambda^\top \bm\Sigma^{-1}$, the expectatio becomes

$$
\bm\Sigma \bm K^\top - \bm\Lambda\bm\Lambda^\top \bm K^\top = \bm\Sigma \bm K^\top - (\bm\Sigma - \bm\phi) \bm K^\top = \bm\phi\bm K^\top
$$
which is not $0$, proving that the imputed scores (by the MAP) are not 0 in expectation and therefore the estimator is not consistant. Bad, very bad.



\section{Asymptotic Properties}


\subsection{Paramatric}

Firs we assume that the parametric space is restrained such that there is a unique solution, and the gllvm is uniquely identified.

First, prove consistency.\\
We can do the asymptotic analysis with Equation 9, where Z is known (since we can approximate).
This thus corresponds to the GEE of a GLM with misspecified variance (it's misspecified since we plug in the map and note the true theta).
We need to show that the M-estimator is well identified (so we need conditions on the weights). According to Nelder (chapter on quasilikelihood), it is consistent even if it is mispecified, as long as the expectation is well specified.
READ GEE BOOK (ZIEGLER).

Second, as thetahat is consistent, the estimator of z is consistent too.
Then we can show that the weights are consistent as well.
So asymptotically we get close the MLE.

\subsubsection{Consistency}

Recall the equation for the Z-estimator:
\begin{equation}
\tilde{\bm\Psi}_n(\bt, \z) = \frac{1}{n} \sum_{i=1}^n\sum_{j=1}^p \bm W_j(\z_{\bt,i}, \bt) y_{ij}- \E_\bt(\bm W_j(\z_{\bt,i}, \bt) y_{ij}|\z_i) ).
\end{equation}
where $\z$ is known and $\z_\bt$ are the ``imputed'' value which depends on $\bt$ and $\y$.
We assume that there is a function $\omega_j$ such that $\bm W_j(\z_{\bt,i}, \bt) \rightarrow \omega_j(\theta)$ as $p\rightarrow \infty$. $\omega$ does not depend on $\y$.
Therefore, we have 
\begin{align}
\tilde{\bm\Psi}_n(\bt, \z) \approx \frac{1}{n} \sum_{i=1}^n\sum_{j=1}^p \omega_j(\bt) y_{ij}- \E_\bt(\omega_j(\bt) y_{ij}|\z_i) ), \\
\approx \frac{1}{n} \sum_{i=1}^n\sum_{j=1}^p \omega_j(\bt) (y_{ij}- \E_\bt(y_{ij}|\z_i) ).
\end{align}
Assumption (uniqueness/identifiability):
\begin{equation}\label{eq. assumption uniqueness}
\E[ \omega_j(\bt) (y_{ij}- \E_\bt(y_{ij}|\z_i) )] = 0  
\end{equation}
has a unique solution at $\bt = \bt_0$.
Assumption \ref{eq. assumption uniqueness} is realistic (we should explain why).
But if the $\z_\bt \rightarrow Z$, and the variance function is not misspecified, we have the equation of the maximum likelihood of the GLM (with known $\z$).
\jb[inline]{Can we infer a bit more? Would be cool if we show under which condition on $\omega_j$ the solution is unique. First, of course we need $\omega\neq 0$... If the model is factor analysis, we would have $\E(y|\z) = \z\bt$ and $\omega(\bt) = \bt$. Thus we would need that $\lim_{p\rightarrow \infty} \E[ \z_\bt \z]$ is nonsingular.
This is similar to linear regression, where one requires $\E[ \z \z]$ nonsingular, but the first $\z$ is replaces by $\z_\bt$.
What do you think?
}




\subsection{Nonparametric}

\section{Computation}



\subsection{Encoder using a NN}


Similarly to variational auto-encoders (VAE) \cite{kingma2014autoencoding}, our estimation method we will use an encoder $h: \mathbb R^p \times \mathbb R^q$, which learns a lower-dimensional representation of the data, and a decoder given by $p_\bt(\x|\z)$, which decodes that representation back into the original data format. Our method uses simulated data from the marginal model $p_\bt(\x)$ to learn the encoder parameters $\bm\xi$. One can then encode $\x$ as $\hat\z = h(\x, \bm\xi)$. With the encoded values substituting for the unknown latent variable, we then use a simple loss function to learn the parameters of the decoder. While we are also interested in the imputing values, our method focuses on the good estimation of the parameters of the decoder, which is justified by an M-estimation argument. 

\subsection{Stochastic-Approximation Algorithm and Autograd}
\section{Applications}
\subsection{GLMM or Missing Variable}
\subsection{GLLVM: SNP}
\subsection{MNIST}
\section{Conclusion}
\newpage

\section{Graveyard}


\subsection{Estimating Function, new presentation}
We consider the class of Generalized Estimating Equations (GEE), which can be written as:

\begin{equation}
\bm\Psi(\y, \z, \bt) = \bm U(\bt, \z; \y) - \E\left[\bm U (\bt, \z, \Y)|\z\right]
\end{equation}

where

\begin{equation}
\bm U(\bt, \z;\y) = \frac{\partial\mu}{\partial\bt}\bm V(\bm\mu)^{-1} \y
\end{equation}

and $\mu = \E[\y|\z]$ and $\V$ is a known variance function . 


If $\z$ is known, these have good properties because because they are 0 in expectation, i.e. they satisfy

$$
\E[\bm\Psi(\y, \z, \bt)|\Z] = 0.
$$

That is, the estimating equation satisfies the first Bartlett identity, one of the fundamental property for consistent estimation. We want to generalize this to the case where $\Z$ is unknown: the $\Z$ need to be marginalized out by their conditional distribution $\Z|\Y$. The estimating equation becomes

\begin{equation}
\bm \Psi(\bt;\y) = \int_{\mathcal Z}\bm\Psi(\y, \z, \bt)f_\bt(\z|\y)d\z.
\label{eqn:gee_marginal}
\end{equation}
These still have good properties, for the same reason as before, because their expectation is 0 (now on the marginal model):

\begin{multline}
\E[\bm\Psi(\y, \bt)] = \int_{\mathcal Y}\int_{\mathcal Z}\bm\Psi(\y, \z, \bt)f_\bt(\z|\y)d\z f_\bt(\y)d\y 
\\= \int_{\mathcal Y}\int_{\mathcal Z}\bm\Psi(\y, \z, \bt)f_\bt(\y|\z)d\y \pi(\z) d\z = 0.
\end{multline}

However maximizing \eqref{eqn:gee_marginal} w.r.t $\bt$ is extremely hard (this would be the EM algorithm, in fact). A way forward would be to use variational inference, or other approximations of $\bm\Psi(\bt, \z, \y)$. Here, we propose the AI trick to approximate \eqref{eqn:gee_marginal}:
%
\begin{equation}
    \tilde{\bm\Psi}(\bt;\y) = \bm U(\bt, \e(\bt, \y); \y) - \E[\bm U(\bt, \e(\bt, \Y); \Y)]
\end{equation}
%
This approximation is good (to be shown ??), but the key here is that, unlike approximations using variational inference or gaussian quadrature, our estimator retains its tractability because its expectation is still 0:

\begin{equation}
    \E[\tilde{\bm\Psi}(\bt;\Y)] = \bm 0.
\end{equation}

\subsection{Some other remarks}

Notice that at the cost of notational complexity, the distribution of $\Y_i$, and indeed of $\Y_i|\Z_i$, can be made conditional on some covariates $\X_i$.
The variance functions $\{V_j\}_{j=1}^p$ in \eqref{eqn:estimator}, similarly to GEEs, can capture dependence among the observations in longitudinal studies, such as when there are multiple measurements for each participant. Alternatively, random effects can be added as latent variables and similarly imputed by the encoder. Our estimator remains consistent by definition and covers all these cases.


\subsection{Summary}
Our estimator is defined as the root $\hat\bt\in\bm\Theta$ of the following so-called estimating equations:
%
\begin{equation}
    \frac{1}{n}\sum_{i=1}^n \bm G(\Y_i, \bt) = \frac{1}{n}\sum_{i=1}^n \E[\bm G(\Y_i, \bt)],
\end{equation}

%
where the expectation is taken under the model $F_{\bt}$. This estimator is a special case of a mimial distance estimator \gb{cite}, which is a large class of estimators which includue method of moments and M-estimators. 

It belongs to the class of M-estimators, which relate an empirical mean of a quantity $\bm G(\Y_i, \bt)$ to its theoretical counterpart\gb{cite something}; it is not unlike the method of moments, except that we allow the quantity to also depend on the parameter. For our purpose, we define this quantity to be
%
\begin{equation}
    \bm G(\Y_i, \bt) = \sum_{j=1}^p \frac{\bm D_j(\Z_i, \bt) Y_{ij}}{V_j(\Z_i, \bt)\phi_j}\Bigg\vert_{\Z_i=\e(\Y_i, \bt)},
\end{equation}
%
where $\mu_{ij}(\Z_i, \bt) = \E[Y_{ij}|\Z_i]$, $\bm D_j(\Z_i, \bt) = \frac{\partial \mu_{ij}(\Z_i, \bt)}{\partial\bt} $, $V_j(\Z_i, \bt)\phi_j = \text{Var}(Y_{ij}|\Z_i)$, $\phi_j$ is a scale parameter and $\e(\Y_i, \bt)$ is the \textit{encoder}, an imputing function estimating the unknown $\Z_i$. By this imputation, the unknown $\Z_i$ are replaced solely by a function of $\Y_i$ and $\bt$, effectively profiling them out of the estimating equations. Note that the quantity $\bm D_j(\Z_i, \bt)$ is computed as if $\Z_i$ were observed. Finally, we allow the distribution of $\Y_i$ to depend on $i$ through external covariates.


\subsection{Special Case of Canonical Parameters}
When $g$ is the cannonical link function, the following relationship holds:
the equation  \eqref{eqn:generalized-linear-models} simplifies, thanks to the relationship

$$
\frac{\partial \mu_i(\bt)}{\partial\bt} = \frac{\partial g^{-1}(\eta_i(\bt))}{\partial \eta_i(\bt)}\frac{\partial \eta_i(\bt)}{\partial\bt} = V_i(\bt)\frac{\partial \eta_i(\bt)}{\partial\bt},
$$

to :

\begin{equation}
    \sum_{i=1}^n \frac{\partial\eta_i(\bt)}{\partial\bt}\frac{(Y_i - \mu_i(\bt))}{\phi}\cdot
    \label{eqn:generalized-linear-models-cannonical}
\end{equation}


{\noindent \em Remainder omitted in this sample. See http://www.jmlr.org/papers/ for full paper.}

% Acknowledgements should go at the end, before appendices and references

\acks{I did that for fun. Please help.}

% Manual newpage inserted to improve layout of sample file - not
% needed in general before appendices/bibliography.

\newpage


\subsection{Derivation as EM?}


\subsubsection{New Draft}
For an exponential family of distribution, the E-step is

$$Q(\theta \mid \theta^{(t)}) = \mathbb{E}_{Z\sim q(z \mid y; \theta^{(t)})} \left[ T(y) {\eta(Z)} - A(\eta(Z)) \right]$$

$$
Q(\theta \mid \theta^{(t)}) = T(y) \, \mathbb{E}_{Z\sim q(z \mid y; \theta^{(t)})} \left[ {\eta(Z, \theta)} \right] - \mathbb{E}_{Z\sim q(z \mid y; \theta^{(t)})} \left[ A(\eta(Z, \theta)) \right]
$$

This is the usual form. It is hard to estimate because $q(z|y; \theta^{(t)})$ is hard etc etc. We propose the following, transformation of the E step (\textcolor{softgreen}{new idea} \textcolor{red}{old idea}):


% $$
% Q(\theta \mid \theta^{(t)}) = T(y) \, \mathbb{E}_{Z\sim q(z \mid y; \theta^{(t)})} \left[ {\eta} \right] - \mathbb{E}_{Y\sim f(y;\theta^{(t)})} \left[ \mathbb{E}_{Z\sim q(Z \mid Y; \theta^{(t)})} \left[ A(\eta(Z)) \right] \right]
% $$

\textcolor{softgreen}{
$$
Q(\theta \mid \theta^{(t)}) = T(y) \, \mathbb{E}_{Z\sim q(z \mid y; \theta^{(t)})} \left[ {\eta(Z, \theta)} \right] - \mathbb{E}_{Y\sim f(y;\theta^{(t)})} \left[ \mathbb{E}_{Z \sim q(z \mid Y; \theta^{(t)})} \left[ A(\eta(Z, \theta)) \right] \right]
$$
}



\textcolor{red}{
$$
Q(\theta \mid \theta^{(t)}) = T(y) \, \mathbb{E}_{Z\sim q(z \mid y; \theta^{(t)})} \left[ {\eta(Z, \theta)} \right] - \mathbb{E}_{Y\sim f(y;\theta^{(t)})} \left[ T(Y)\mathbb{E}_{Z\sim q(z \mid Y; \theta^{(t)})} \left[ \eta(Z, \theta) \right] \right]
$$
}




The M step is

$$
\theta^{(t+1)} = \arg \max_\theta \, Q(\theta \mid \theta^{(t)}).
$$

Artenatively, generalized modified EM or stochastica approximation EM can be used.


\subsubsection{Gradient computation}
The gradient is

$$
\frac{\partial Q(\theta \mid \theta^{(t)})}{\partial \theta} = T(y) \, \mathbb{E}_{Z \sim q(z \mid y; \theta^{(t)})} \left[ \frac{\partial \eta(Z, \theta)}{\partial \theta} \right] - \mathbb{E}_{Y \sim f(y; \theta^{(t)})} \left[ T(Y) \, \mathbb{E}_{Z \sim q(z \mid Y; \theta^{(t)})} \left[ \frac{\partial \eta(Z, \theta)}{\partial \theta} \right] \right].
$$



\subsubsection{Fisher Consistency}
The solution to the EM algorithm (assumed unique...), denoted as $\hat\theta$ satisfies:

$$
Q(\theta|\hat\theta)|_{\theta=\hat\theta} = 0
$$
%
Thus, the estimator can alternatively be obtained as the root of the following estimating equations
%
$$
T(y) \, \mathbb{E}_{q(z \mid y; \theta)} \left[ \frac{\partial {\eta(Z, 
 \theta)}}{\partial \theta} \right] - \mathbb{E}_{Y \sim f(y; \theta)} \left[ T(Y) \cdot \mathbb{E}_{q(z \mid Y; \theta)} \left[ \frac{\partial {\eta(Z, \theta)}}{\partial \theta} \right] \right] = 0,
$$
%
which is satisfied under the  true model, thus insuring Fisher consistency. Notice that $q$ needs not be exact, it can be approximate, to be defined later.


\subsubsection{Show:}

\begin{itemize}
    \item Asymptotic Equivalence to the MLE under the true posterior function $q$ (easy to show by definition of conditional distribution...??!!?? i think so, see below).
    \item Robustness to misspecification of the $q$ function... can we quantify how far can we be?
    \item Does the algo converge??!! ??!!?!?! (in practice yes, but under which conditions)
\end{itemize}

Given n  independent observations  $(Y_i, Z_i) \sim f(y, z; \theta)$, we have:
$$
\frac{1}{n} \sum_{i=1}^n \mathbb{E}_{Z_i \mid Y_i} \left[ h(Y_i, Z_i, \theta) \right] = \mathbb{E}_{Y} \left[ \mathbb{E}_{Z \mid Y} \left[ h(Y, Z, \theta) \right] \right] + o\left(\frac{1}{n}\right).
$$




\begin{align}
\sum_{i=1}^n T(y_i) \, \mathbb{E}_{Z_i\sim q(z \mid y_i; \theta)} \left[ \frac{\partial {\eta(Z_i, 
 \theta)}}{\partial \theta} \right] - \sum_{i=1}^n   \mathbb{E}_{Z_i \sim q(z \mid y_i; \theta)} \left[ \mathbb{E}[T(Y)|\eta(Z_i, \theta)] \,\frac{\partial \eta(Z_i, \theta)}{\partial \theta} \right] = 0 \\
 \sum_{i=1}^n T(y_i) \, \mathbb{E}_{Z_i\sim q(z \mid y_i; \theta)} \left[ \frac{\partial {\eta(Z_i, 
 \theta)}}{\partial \theta} \right] - \sum_{i=1}^n   \mathbb{E}_{Z_i \sim q(z \mid y_i; \theta)} \left[ \mathbb{E}[T(Y)|\eta(Z_i, \theta)] \,\frac{\partial \eta(Z_i, \theta)}{\partial \theta} \right] = 0
\end{align}




\newpage

\subsubsection{Benji's derivation}
Assume that $Z, Y$ is a couple of random variables with joint distribution of the form 
$$ f_{Y,Z\mid \theta}(y,z) = h(y) \exp\left( \eta_{\theta}(z) T(y) - A(\eta_{\theta}(z))\right) $$
where $\theta \in \R^{qqchose}$, for some (nice) functions $h$, $\eta$ and $A$.

In what follows, expectations will be written with a subscript indicating the value of $\theta$ taken for the distribution of $(X, Y)$, $f_{Y,Z\mid \theta}$. That way, for any function $g$, $$E_\theta\left[g(Y, Z)\right] = \int g(y, z) f_{Y, Z \mid \theta}(y, z)dy dz \,.$$

The first Bartlett identity yields
$$\fracpartial{A}{\eta} \left(\eta_\theta(z)\right) \fracpartial{\eta_\theta}{\theta}(z) = \E_{\theta}\left[T(Y) \mid Z = z\right]\, .$$

Given some observed $y$, call $\ell(\theta)$ the log-likelihood of the observations under the model:
$$\ell(\theta) = \E_{\theta}[ \log(f_{Y, Z\mid \theta}(Y, Z) \mid Y = y ] = \E_{\theta}[ T(Y) \eta_{\theta}(Z) - A(\eta_{\theta}(Z)) \mid Y = y] \, .$$
The classical E-M algorithm optimizes $\ell(\te)$ by defining 
$$Q(\theta_1, \theta_2)  = \E_{\theta_2}\left[ T(Y) \, \eta_{\theta_1}(Z) - A(\eta_{\theta_1}(Z)) \mid Y = y\right] \, ,$$
such that $\ell(\theta) = Q(\theta, \theta)$, 
and by repeating the following two steps: 
\begin{itemize}
    \item E-step: Given the current parameter estimate $\te^{(t)}$, compute the function $\te_1 \mapsto Q(\te_1|\te^{(t)})$.
    \item M-step: Maximize this function with respect to $\te_1$:
    $$\te^{(t+1)} = \textrm{argmax}_{\te_1} Q(\te_1|\te^{(t)}) \, .$$
\end{itemize}
The resulting sequence $\te^{(1)}, \te^{(2)}, \ldots, $ is guarantied to be such that $\ell(\te^{(1)}), \ell(\te^{(2)}), \dots$ is increasing.

We can rewrite $Q(\theta_1, \theta_2)$ as 
$$Q(\theta_1, \theta_2) = T(y) \E_{\theta_2}\left[\eta_{\theta_1}(Z)\mid Y = y \right] - \E_{\theta_2}\left[A(\eta_{\theta_1}(Z)) \mid Y = y\right] .$$

Our estimator can be defined by replacing, in this expression of $Q(\theta_1, \theta_2)$, the conditional expectation $\E_{\theta_2}\left[A(\eta_{\theta_1}(Z)) \mid Y = y \right]$ by the total expectation $\E_{\theta_2}\left[A(\eta_{\theta_1}(Z))\right]$.

Call $\tilde Q(\theta_1, \theta_2)$ the version of $Q(\theta_1, \theta_2)$ modified in this manner:
$$\tilde {Q}(\theta_1, \theta_2) = T(y) \E_{\theta_2}\left[\eta_{\theta_1}(Z)\mid Y = y \right] - \E_{\theta_2}\left[A(\eta_{\theta_1}(Z))\right] \, . $$

Our estimator is then the $\te$ value towards which converges a modified E-M algorithm, in which the function ${Q}(\theta_1, \theta_2)$ is replaced by the function $\tilde {Q}(\theta_1, \theta_2)$.

To see that more clearly, consider the following derivation. 

When convergence of these this modified E-M algorithm is met, the corresponding value $\theta$ must satisfy
 $\fracpartial{\tilde{Q}}{\te_1}(\te_1 = \te,\te_2 = \te) = 0$.
Under suitable regularity conditions, 
$$\fracpartial{\tilde{Q}}{\te_1}(\te_1, \te_2) = T(y) \E_{\te_2}\left[ \fracpartial{\eta_{\te_1}}{\te_1}(Z)\mid Y = y \right] - \E_{\te_2} \left[ \fracpartial{A}{\eta}(\eta_{\te_1}(Z)) \,\fracpartial{\eta_{\te_1}}{\te_1}(Z) \right] \,.$$ 
As $\fracpartial{A}{\eta}(\eta_{\te_1}(Z)) = \E_{\te_2}\left[T(Y)\mid Z\right]$, 

$$\E_{\te_2}\left[T(Y)\mid Z\right] \,\fracpartial{\eta_{\te_1}}{\te_1}(Z) = \E_{\te_2}\left[T(Y) \,\fracpartial{\eta_{\te_1}}{\te_1}(Z) \mid Z\right]$$
and therefore

$$\E_{\te_2} \left[ \fracpartial{A}{\eta}(\eta_{\te_1}(Z)) \,\fracpartial{\eta_{\te_1}}{\te_1}(Z) \right] = \E_{\te_2}\left[ \E_{\te_2}\left[T(Y) \,\fracpartial{\eta_{\te_1}}{\te_1}(Z) \mid Z\right] \right] = \E_{\te_2}\left[T(Y) \,\fracpartial{\eta_{\te_1}}{\te_1}(Z) \right] \, .$$

At convergence at, we then have that 
$$ T(y) \, \E_\theta\left[ \fracpartial{\eta_\theta}{\theta}(Z) \bigm\vert Y = y \right] = \E_\theta\left[ T(Y) \, \fracpartial{\eta_\theta}{\theta}(Z) \right] \, .$$

This can also be written less concisely as 
$$ T(y) \, \E_\theta\left[ \fracpartial{\eta_\theta}{\theta}(Z) \bigm\vert Y = y \right] = \E_\theta\left[ T(Y) \, \E_\theta\left[ \fracpartial{\eta_\theta}{\theta}(Z) \mid Y\right] \right] \, $$
to stress that this estimating equation corresponds to a method of moment estimator, or to a Fisher-consistent M-estimator.

We generalize this estimating equation by allowing the conditional expectations in both side of the equation to be replaced by some approximation -- the same approximation must be used in both sides. Calling such an approximation $\tilde{E_\theta}(\fracpartial{\eta_\theta}{\theta}(Z), y)$, we obtain the following estimating equation:

$$ T(y) \, \tilde{E}_\theta\left[ \fracpartial{\eta_\theta}{\theta}(Z) , y \right] = \E_\theta\left[ T(Y) \, \tilde{E}_\theta\left[ \fracpartial{\eta_\theta}{\theta}(Z) , Y\right] \right] \, .$$


Typically one can use $\tilde{E_\theta}(\fracpartial{\eta_\theta}{\theta}(Z), y) = \fracpartial{\eta_\theta}{\theta}(z_{MAP})$ where $z_{MAP} = \textrm{argmax}_z f_{Z \mid Y}(z)$ or




















\newpage

\textcolor{red}{
Call $g_{\theta, \theta^{(t)}}(y) = \mathbb{E}_{q(z \mid y, \tz; \theta^{(t)})} \left[ A(\eta_\theta(z)) \right]$
$$
Q(\theta \mid \theta^{(t)}) = T(y) \, \mathbb{E}_{Z\sim q(z \mid y; \theta^{(t)})} \left[ {\eta} \right] - \mathbb{E} \left[ g_{\theta^{(t)}}(Y) \right]
$$
}

\newpage
\appendix
\section*{Appendix A.}
\label{app:theorem}

% Note: in this sample, the section number is hard-coded in. Following
% proper LaTeX conventions, it should properly be coded as a reference:

%In this appendix we prove the following theorem from
%Section~\ref{sec:textree-generalization}:

In this appendix we prove the following theorem from
Section~6.2:

\noindent
{\bf Theorem} {\it Let $u,v,w$ be discrete variables such that $v, w$ do
not co-occur with $u$ (i.e., $u\neq0\;\Rightarrow \;v=w=0$ in a given
dataset $\dataset$). Let $N_{v0},N_{w0}$ be the number of data points for
which $v=0, w=0$ respectively, and let $I_{uv},I_{uw}$ be the
respective empirical mutual information values based on the sample
$\dataset$. Then
\[
	N_{v0} \;>\; N_{w0}\;\;\Rightarrow\;\;I_{uv} \;\leq\;I_{uw}
\]
with equality only if $u$ is identically 0.} \hfill\BlackBox

\noindent
{\bf Proof}. We use the notation:
\[
P_v(i) \;=\;\frac{N_v^i}{N},\;\;\;i \neq 0;\;\;\;
P_{v0}\;\equiv\;P_v(0)\; = \;1 - \sum_{i\neq 0}P_v(i).
\]
These values represent the (empirical) probabilities of $v$
taking value $i\neq 0$ and 0 respectively.  Entropies will be denoted
by $H$. We aim to show that $\fracpartial{I_{uv}}{P_{v0}} < 0$....\\

{\noindent \em Remainder omitted in this sample. See http://www.jmlr.org/papers/ for full paper.}



\vskip 0.2in
\bibliography{references}

\end{document}